% !TEX program = xelatex
% 三亚学院本科论文参考工作稿；未取得学院备案版式，不能声称官方复刻。
\documentclass[12pt,a4paper,fontset=none]{ctexart}
% 以下默认参数均待所在学院备案模板核对，不是已取得的官方规范。
\usepackage[top=30mm,bottom=25mm,left=30mm,right=20mm,headheight=15pt]{geometry}
% MIT; default preview uses TeX-distributed fonts, NOT school typography.
\RequirePackage{fontspec,xeCJK,etoolbox}
\providecommand{\ThesisFontMode}{preview}
\ifdefstring{\ThesisFontMode}{school}{%
  \newcommand{\ThesisCheckFont}[1]{\IfFontExistsTF{#1}{}{\PackageError{thesis-fonts}{Missing font: #1}{Install legally licensed fonts or use THESIS_FONTS=preview.}}}
  \ThesisCheckFont{Times New Roman}\ThesisCheckFont{SimSun}
  \ThesisCheckFont{SimHei}\ThesisCheckFont{FangSong}\ThesisCheckFont{KaiTi}
  \setmainfont{Times New Roman}
  \setCJKmainfont[AutoFakeBold=2]{SimSun}
  \setCJKsansfont[AutoFakeBold=2]{SimHei}
  \setCJKmonofont{FangSong}
  \setCJKfamilyfont{zhsong}[AutoFakeBold=2]{SimSun}
  \setCJKfamilyfont{zhhei}[AutoFakeBold=2]{SimHei}
  \setCJKfamilyfont{zhkai}[AutoFakeBold=2]{KaiTi}
  \setCJKfamilyfont{zhfs}{FangSong}
}{%
  \ifdefstring{\ThesisFontMode}{preview}{}{\PackageError{thesis-fonts}{Invalid font mode}{Use preview or school.}}
  \setmainfont{texgyretermes-regular.otf}[BoldFont=texgyretermes-bold.otf,ItalicFont=texgyretermes-italic.otf,BoldItalicFont=texgyretermes-bolditalic.otf]
  \setCJKmainfont{FandolSong-Regular.otf}[BoldFont=FandolSong-Bold.otf]
  \setCJKsansfont{FandolHei-Regular.otf}[BoldFont=FandolHei-Bold.otf]
  \setCJKmonofont{FandolFang-Regular.otf}
  \setCJKfamilyfont{zhsong}{FandolSong-Regular.otf}[BoldFont=FandolSong-Bold.otf]
  \setCJKfamilyfont{zhhei}{FandolHei-Regular.otf}[BoldFont=FandolHei-Bold.otf]
  \setCJKfamilyfont{zhkai}{FandolKai-Regular.otf}[AutoFakeBold=2]
  \setCJKfamilyfont{zhfs}{FandolFang-Regular.otf}
}
% Geometric/shapes may not exist in the serif font; ctex's open symbol font.
\xeCJKDeclareCharClass{CJK}{"2460 -> "24FF,"3200 -> "32FF}

\providecommand{\songti}{\CJKfamily{zhsong}}
\providecommand{\heiti}{\CJKfamily{zhhei}}
\usepackage{amsmath,amssymb,graphicx,booktabs,tabularx,longtable,caption}
\usepackage{titlesec,tocloft,fancyhdr,setspace,indentfirst,listings}
\usepackage{tikz}\usetikzlibrary{positioning,arrows.meta}
\usepackage[super]{gbt7714}
\usepackage[hidelinks]{hyperref}
\setcounter{tocdepth}{3}\setcounter{secnumdepth}{3}
\titleformat{\section}{\heiti\zihao{-3}}{\thesection}{1em}{}
\titleformat{\subsection}{\heiti\zihao{4}}{\thesubsection}{1em}{}
\titleformat{\subsubsection}{\heiti\zihao{-4}}{\thesubsubsection}{1em}{}
\renewcommand{\cfttoctitlefont}{\hfill\heiti\zihao{-3}}
\renewcommand{\cftaftertoctitle}{\hfill\mbox{}}
\tocloftpagestyle{sanya}
\fancypagestyle{sanya}{\fancyhf{}\fancyhead[L]{\songti\zihao{5}三亚学院毕业论文（设计）}\fancyhead[R]{\songti\zihao{5}第\thepage 页}\renewcommand{\headrulewidth}{0.4pt}}
\DeclareCaptionFont{sycap}{\heiti\zihao{5}}
\captionsetup{font=sycap,labelsep=space,justification=centering}
\captionsetup[table]{position=top}
\counterwithin{figure}{section}\counterwithin{table}{section}\counterwithin{equation}{section}
\lstset{basicstyle=\ttfamily\small,breaklines=true,columns=fullflexible,frame=single,numbers=left,numberstyle=\tiny,showstringspaces=false}

\newcommand{\thesistitle}{应用型研究论文写作与排版示例}
\newcommand{\thesistitleeng}{A Writing and Typesetting Example for an Applied Research Thesis}
\newcommand{\thesisauthor}{张三}
\newcommand{\thesisid}{2022000001}
\newcommand{\thesiscollege}{请填写所在学院}
\newcommand{\thesismajor}{请填写专业与班级}
\newcommand{\thesisadvisor}{请填写校内指导教师}
\newcommand{\thesisindustryadvisor}{请填写行业导师（如适用）}
\begin{document}
\songti\zihao{-4}\setstretch{1.5}\setlength{\parindent}{2em}
\begin{titlepage}\thispagestyle{empty}\centering
\vspace*{10mm}{\heiti\zihao{1}三亚学院}\par\vspace{15mm}
{\heiti\zihao{2}本科毕业论文（设计）}\par\vspace{15mm}
{\heiti\zihao{-2}\thesistitle}\par\vspace{18mm}
\renewcommand{\arraystretch}{1.8}
\begin{tabular}{r>{\raggedright\arraybackslash}p{.65\linewidth}}
姓名：&\thesisauthor\\学号：&\thesisid\\
学院：&\thesiscollege\\专业班级：&\thesismajor\\
校内导师：&\thesisadvisor\\行业导师：&\thesisindustryadvisor
\end{tabular}\vfill
二〇二六年六月\par\medskip
{\heiti\zihao{5}非官方参考工作稿——封面及版式待学院备案模板核对}
\end{titlepage}
\pagestyle{empty}\section*{独创性声明（占位）}
请在此插入学院当年要求的独创性声明原文。当前未取得统一声明表，故不编造“官方声明”。提交前应由本人确认论文为独立完成的真实成果，注明引用和协助，并按要求亲笔签名。不得保留此提示作为正式声明。
\vspace{25mm}\begin{flushright}作者签名：\underline{\makebox[30mm]{}}\par\medskip
日期：\underline{\makebox[40mm]{}}\end{flushright}
\clearpage\pagenumbering{Roman}\pagestyle{sanya}
\phantomsection\addcontentsline{toc}{section}{摘要}
\begin{center}{\heiti\zihao{-2}\thesistitle}\par\medskip{\heiti\zihao{-3}摘\quad 要}\end{center}
本文为本科论文写作与排版工作稿，不报告真实研究成果。模板以应用型研究为线索，组织问题定义、文献综述、方法设计、证据报告、讨论及结论，提供流程图、三线表、数学公式、程序清单和附录等常用元素。文中不预填调查样本量、企业数据或实验性能，避免把排版示例误认为实证发现。三亚学院2026届工作通知明确要求学院结合学科特点确定格式并报教务处审核备案，因此本文件的页面、字号、封面和页码仅为可修改起点，不是已核验的学院模板。正式写作应取得所属学院备案版式，依据实际完成的研究替换全部示例内容，并核对签名、检测、电子文档及归档要求。
\par\medskip\noindent\textbf{关键词：}应用研究；研究设计；论文写作；排版示例
\clearpage\phantomsection\addcontentsline{toc}{section}{Abstract}
\begin{center}{\bfseries\zihao{-2}\thesistitleeng}\par\medskip{\bfseries\zihao{-3}Abstract}\end{center}
This document is an unofficial working template rather than a report of empirical findings. It demonstrates problem definition, literature review, methods, evidence reporting, discussion, and conclusions, together with diagrams, tables, equations, code, and appendices. No survey, experiment, or company investigation was conducted for this example. The university's notice for the 2026 cohort requires colleges to establish discipline-specific formats and submit them for approval. Consequently, the present layout must be checked against the relevant college's approved template before submission.
\par\medskip\noindent\textbf{Keywords:} applied research; research design; thesis writing; typesetting
\clearpage\tableofcontents
\clearpage\pagenumbering{arabic}
\section{绪论}
\subsection{问题定义与应用情境}
应用型研究需要说明真实情境中的问题、利益相关方和约束。应把“研究对象是什么”“希望回答什么问题”“可以取得哪些证据”分开写清楚。选题不能只在题名中添加企业或城市名称而缺少相应资料，也不应预设必须得到正面的结果。
\subsection{文献综述与已有方法}
阅读原始资料，比较方法的假设、适用条件和证据质量，而不是将若干摘要拼接成综述。下面引用真实出版物，仅用于检验中文和英文引文排版\cite{zhou2016,lamport1994}；实际论文须根据选题重新检索和核对。
\subsubsection{概念及分析单位}
为关键概念给出操作性定义，说明分析单位是个体、事件、产品还是组织。样本数、记录数和重复测量次数不应混用。本脚注演示解释性补充\footnote{本文件没有收集企业或个人资料，所有研究场景均为写作提示。}。
\subsection{研究内容与章节安排}
章节安排应服务于问题和证据链。工科研究可围绕需求、设计、实现和测试组织；文经管类可围绕概念、资料、分析和建议组织。学校通知要求各学院按学科特点确定格式，不能把本工作稿的章节清单当成统一要求。
\clearpage\section{研究设计与实施方案}
\subsection{资料来源与使用权限}
报告来源、版本、时间范围和取得方式。企业内部材料应取得适当授权，个人资料应按要求匿名化。无法公开原始材料时，应说明限制，并提供允许披露的变量定义和处理步骤，不能因资料受限而编造数字。
\subsection{技术路线}
图~\ref{fig:route} 展示可修改的技术路线。每一步均应有明确产物，例如研究问题表、资料清单、分析脚本和结果核查记录。
\begin{figure}[htbp]\centering
\begin{tikzpicture}[node distance=6mm,b/.style={draw,minimum width=72mm,minimum height=10mm,font=\small},>=Stealth]
\node[b](a){明确问题、边界与评价准则};
\node[b,below=of a](b){获得资料与必要授权};
\node[b,below=of b](c){实施研究并保留过程记录};
\node[b,below=of c](d){分析证据、复核结果与讨论局限};
\draw[->](a)--(b);\draw[->](b)--(c);\draw[->](c)--(d);
\end{tikzpicture}
\caption{研究组织路线示例}\label{fig:route}
\end{figure}
\subsection{评价指标与数学表达}
指标应与研究目标对应，给出分母、单位和计算规则。例如平均绝对误差可写为：
\begin{equation}\label{eq:mae}\operatorname{MAE}=\frac{1}{n}\sum_{i=1}^{n}|y_i-x_i|.\end{equation}
式中，$x_i$ 为参考量，$y_i$ 为观测量，$n$ 为有效观测数。并非每种应用研究都应使用这一指标，不能为了插入公式而选用不适合的方法。
\subsection{复现与质量控制}
记录软件版本、参数、资料处理规则和异常情况。适用时应使用独立验证、交叉核查或反例分析，区分质量控制与事后挑选结果。研究涉及人类参与者、动物或敏感数据时，须遵守相应伦理与安全要求。
\clearpage\section{结果报告与分析}
\subsection{结果呈现方式}
真实论文应先说明资料规模、缺失和排除，再报告与研究问题直接相关的结果。表~\ref{tab:evidence} 是证据追踪表的结构示例，不包含真实调查结论。
\begin{table}[htbp]\centering\caption{证据追踪表结构示例}\label{tab:evidence}
\begin{tabularx}{\linewidth}{lXX}\toprule 研究问题 & 应保留的证据 & 分析限制\\\midrule
问题一 & 来源记录、原始变量与处理规则 & 样本覆盖与来源偏差\\
问题二 & 方案说明、对照条件与测试记录 & 混杂因素和评价范围\\
问题三 & 分析过程、异常记录与反例 & 外推条件与不确定性\\\bottomrule\end{tabularx}
\end{table}
\subsection{事实、解释与推测}
把“资料显示了什么”“作者如何解释”和“尚需验证的推测”分开。相关性不能直接支持因果结论，个别成功案例也不能证明普遍有效。图表必须有来源与正文引导，避免重复堆放同一组数据。
\subsection{负面结果与异常记录}
如实保留没有达到预期的结果，说明处理方法及其影响。不能把未开展的测试写成已完成，也不能以预期目标充当实际结果。本模板没有实证发现，因此不提供任何性能或政策效果结论。
\clearpage\section{讨论、建议与结论}
\subsection{与已有证据的比较}
讨论应将本研究放回相关文献和应用情境，比较方法及条件差异。若提出实践建议，需要说明实施主体、资源、风险及适用前提，不宜只列空泛口号。
\subsection{局限与未来工作}
说明数据、方法和实施条件带来的限制，区分可在本研究中解决的问题和未来研究计划。若未能完成某项工作，应明确记为未完成，不得使用“验证了”“证明了”等措辞。
\subsection{结论}
结论应回答研究问题并与结果一致，不引入新证据。本文件只演示论文结构和版式，不能作为研究成果提交。正式使用前还须取得所在学院的备案格式，核对检测、答辩和归档要求。

\clearpage\begingroup
\renewcommand{\bibsection}{\section*{参考文献}\addcontentsline{toc}{section}{参考文献}}
\renewcommand{\bibfont}{\songti\zihao{5}\setstretch{1}}
\bibliographystyle{gbt7714-2015-numeric}\bibliography{refs}
\endgroup
\clearpage\section*{致\quad 谢}\addcontentsline{toc}{section}{致谢}
请根据真实情况感谢指导教师、行业导师、提供研究条件的单位和给予具体帮助的人员。不应虚构调研经历或他人贡献；涉及企业名称和个人身份时，还需确认披露权限。此处为写作提示，正式论文请替换。
\clearpage\section*{附\quad 录}\addcontentsline{toc}{section}{附录}
\subsection*{附录 A\quad 数据字典与归档清单}
\begin{longtable}{lp{.65\linewidth}}\toprule 条目 & 说明\\\midrule\endhead
变量定义 & 名称、含义、单位、允许值与缺失编码\\
资料来源 & 原始地址、版本、许可和取得日期\\
分析程序 & 输入、输出、运行环境和执行命令\\
过程材料 & 任务书、开题、中期、评阅、答辩等按学院要求整理\\\bottomrule\end{longtable}
\subsection*{附录 B\quad 最小计算示例}
\begin{lstlisting}[language=Python]
# Synthetic values for a typesetting example only.
reference = [1.0, 2.0, 3.0]
observed = [1.1, 1.9, 3.0]
mae = sum(abs(y - x) for x, y in zip(reference, observed)) / len(reference)
print(mae)
\end{lstlisting}
附录顺序与位置以学院规定为准；本文件的示例不替代官方过程管理表格。

\end{document}
